\subsection{48维以上：从块库到像集的并}

最初的构造族采用 Pless 邻格中63个互不相交的块。
这些块的前缀大小与根系尾码结合，给出49至55维的七项构造。
配套论证保留邻格的定义及其最小范数证明。

较大的构造族采用公开目录中的 $P_{48p}$ 基底。
公开的7069条极小直线类被扩为7077条，
随后以像集之并的大小为目标，联合选择它的自同构像。
将重复直线分配给第一次出现的位置，就得到互不相交的子类，
同时保留每个类内部的内积界。

七个并集大小分别是
\[
7077,\quad21231,\quad42459,\quad84874,\quad141328,\quad253851,\quad442507.
\]
母壳与根系尾码之外，每条直线贡献两个点。
50维见证使用三个不相交像；较大的像族有重叠，
其影响已经精确计入并集大小。

不同头类分配到不同的反极尾直线，
从而具有跨类内积界
$(3/4)(1/2)+(1/4)(1/2)=1/2$。
因此，尾标签的单射分配与所选头类的并共同构成有限输入。
